\subsection{Estimand and regression estimators}
Units $i=1,\dots,n$ sit on a fixed undirected graph with adjacency $A$, degrees $d_i$ and row-normalised adjacency $W=D^{-1}A$. A design draws an assignment $\bz\in\{0,1\}^n$; the outcome is $Y_i=\mu_i(\bz)+\varepsilon_i$ with a mean function $\mu_i$ that may depend on all of $\bz$ and mean-zero noise $\varepsilon_i$. The estimand is the global average treatment effect
\[
\gate=\frac1n\sum_{i=1}^n\bigl[\mu_i(\mathbf 1)-\mu_i(\mathbf 0)\bigr].
\]
Since outcomes are synthetic, $\gate$ is computed \emph{exactly} from the noise-free mean function; there is no Monte-Carlo error in the target.

A \emph{regression-based} estimator fits a model $\hat f$ from features $x_i(\bz)$ built from the assignment and the graph to the observed outcomes of one experiment, and returns
\[
\hat\gate=\frac1n\sum_{i=1}^n\bigl[\hat f\bigl(x_i(\mathbf 1)\bigr)-\hat f\bigl(x_i(\mathbf 0)\bigr)\bigr].
\]
The key structural fact is that the two prediction points $x_i(\mathbf 1),x_i(\mathbf 0)$ generally lie outside the support of the features observed under the design: under Bernoulli$(p)$ randomisation the treated-neighbour fraction $f_i=(W\bz)_i$ concentrates around $p$, while the predictions need $f_i=1$ and $f_i=0$. We study the following feature maps (all include the intercept and $z_i$): \texttt{ols\_f} adds $f_i$ (the setting of \citealt{chin2019regression}); \texttt{ols\_cnt} replaces it by the scaled count $(A\bz)_i/\bar d$; \texttt{ols\_f2} adds the two-hop fraction $(W^2\bz)_i$; \texttt{ols\_poly} adds $f_i^2$ and $z_if_i$; \texttt{ols\_f\_deg} adds log-degree and its interactions; \texttt{ridge\_rich} fits a ridge regression on a rich set of such features; and \texttt{gbm} is a histogram gradient-boosted tree regressor on $(z_i,f_i,(W^2\bz)_i,\text{log-degree},\text{clustering})$. Baselines are the difference in means (DM) and a H\'ajek estimator contrasting units whose whole closed neighbourhood is treated versus control, with inverse probabilities of that event under the design \citep{aronow2017estimating,ugander2013graph}.

\subsection{Error decomposition}
For each estimator we separate (i) the \emph{approximation error}, the error of the estimator when it is fit to the noise-free, baseline-free mean function $\mu(\bz)$ on the same assignment (a ``population'' fit), which isolates what the model class can do from the design's exposure support, from (ii) the \emph{estimation error} due to finite $n$, baselines and noise. Across replicates we report bias, standard deviation and RMSE, with RMSE relative to $|\gate|$ as the scale-free primary metric: the question is how much of the error is systematic (does not shrink with $n$) and how much is variance.

\subsection{Analytic predictions in the linear-exposure family}
We derive first-order predictions for the OLS$(z,f)$ estimator under Bernoulli$(p)$ and compare them to simulation (Section~\ref{sec:experiments}). We stress that these are large-$n$ linearisation arguments checked numerically, not finite-sample theorems. Write the mean function as $\mu_i(\bz)=b_i+\delta z_i+g(f_i)$ with baselines $b_i$ and direct effect $\delta$, so that $\gate=\delta+g(1)-g(0)$.
\begin{description}
\item[P1 (nonlinear exposure response).] Because $z_i$ is independent of $f_i$, the population OLS slope on $f$ is $\beta_f^*=\mathrm{Cov}(g(f),f)/\mathrm{Var}(f)$ with $f_i\mid d_i\sim\mathrm{Bin}(d_i,p)/d_i$, and the asymptotic bias of $\hat\gate$ is $\beta_f^*-[g(1)-g(0)]$. It vanishes for linear $g$; for smooth $g$ and high degree $\beta_f^*\to g'(p)$, so the bias tends to a tangent-versus-secant gap that does not shrink with $n$.
\item[P2 (variance, correct specification).] With a fixed centred baseline vector $b$ and noise variance $\sigma_e^2$, $\mathrm{Var}(\hat\gate)\approx\bigl[\|b+d_HW^\top b\|^2/n+\sigma_e^2(1+d_H)\bigr]/\bigl(np(1-p)\bigr)$, where $d_H$ is the harmonic mean degree. The error is design-based (the baselines projected on random exposures contribute), grows with degree, and is independent of the interference strength.
\item[P3 (linear in $\bz$ but not in $f$).] If $\mu_i(\bz)=\delta z_i+\sum_jM_{ij}z_j$ (heterogeneous spillovers, two-hop terms, count-scaled exposure), the OLS limit has $z$-slope $\delta+\mathrm{mean}_iM_{ii}$ and $f$-slope $\sum_{ij}W_{ij}M_{ij}/\sum_{ij}W_{ij}^2$. For $M=\mathrm{diag}(\gamma)W$ the $f$-slope is a $1/d_i$-weighted mean of $\gamma_i$ whereas $\gate$ weights units equally, so the estimator is biased whenever spillover strength correlates with degree.
\end{description}
Derivations are in the supplementary material (\texttt{code/docs/THEORY.md}); their numerical checks are in \texttt{results/theory\_check\_*.csv}.
